%% Final edits by Jussi Kangasharju and Pirjo Moen
%% This file is modified by Veli Mäkinen from HY_fysiikka_LuKtemplate.tex authored by Roope Halonen ja Tomi Vainio.
%% Some text is also inherited from engl_malli.tex by Kutvonen, Erkiö, Mäkelä, Verkamo, Kurhila, and Nykänen.


% STEP 1: Choose oneside or twoside
\documentclass[english,twoside,openright]{UH_DS_MSc}
%finnish,swedish

%\usepackage[utf8]{inputenc} % For UTF8 support. Use UTF8 when saving your file.
\usepackage{lmodern} % Font package
\usepackage{textcomp} % Package for special symbols
\usepackage[pdftex]{color, graphicx} % For pdf output and jpg/png graphics
\usepackage[pdftex, plainpages=false]{hyperref} % For hyperlinks and pdf metadata
\usepackage{fancyhdr} % For nicer page headers
\usepackage{tikz} % For making vector graphics (hard to learn but powerful)
%\usepackage{wrapfig} % For nice text-wrapping figures (use at own discretion)
\usepackage{amsmath, amssymb} % For better math
%\usepackage[square]{natbib} % For bibliography
\usepackage[footnotesize,bf]{caption} % For more control over figure captions
\usepackage{titlesec}
\usepackage[titletoc]{appendix}

% The class sets a headheight fractionally smaller than fancyhdr wants for this
% font size, which otherwise emits a warning on every page.
\setlength{\headheight}{15pt}
\addtolength{\topmargin}{-3pt}

\onehalfspacing %line spacing
%\singlespacing
%\doublespacing

%\fussy 
\sloppy % sloppy and fussy commands can be used to avoid overlong text lines

% STEP 2:
% Set up all the information for the title page and the abstract form.
% Replace parameters with your information.
\title{Performance, Safety, and Engineering Trade-offs of Rust and C++ in Nokia RAN Control-Plane Systems}
\author{Ha Vien}
\date{\today}
\prof{Keijo Heljanko, Christoph Luneburg, Petteri Kotimaa}
\censors{Keijo Heljanko}{}{}
\keywords{Rust, C++, RAN control plane, tail latency, memory safety, fault injection, FFI}
\depositeplace{Helsinki University Library}
\additionalinformation{Software study track}


\classification{\protect{\ \\
\  Software and its engineering $\rightarrow$ Software notations and tools $\rightarrow$ General programming languages\  \\
\  Computer systems organization $\rightarrow$ Real-time systems $\rightarrow$ Real-time system specification\\
\  Networks $\rightarrow$ Network types $\rightarrow$ Mobile networks\\
}}

% if you want to quote someone special. You can comment this line and there will be nothing on the document.
%\quoting{Bachelor's degrees make pretty good placemats if you get them laminated.}{Jeph Jacques} 


% OPTIONAL STEP: Set up properties and metadata for the pdf file that pdfLaTeX makes.
% But you don't really need to do this unless you want to.
\hypersetup{
    %bookmarks=true,         % show bookmarks bar first?
    unicode=true,           % to show non-Latin characters in Acrobat’s bookmarks
    pdftoolbar=true,        % show Acrobat’s toolbar?
    pdfmenubar=true,        % show Acrobat’s menu?
    pdffitwindow=false,     % window fit to page when opened
    pdfstartview={FitH},    % fits the width of the page to the window
    pdftitle={Performance, Safety, and Engineering Trade-offs of Rust and C++ in Nokia RAN Control-Plane Systems},
    pdfauthor={Ha Vien},
    pdfsubject={Master's thesis, Computer Science},
    pdfcreator={Ha Vien},
    pdfproducer={pdfLaTeX}, % producer of the document
    pdfkeywords={Rust} {C++} {RAN control plane} {tail latency} {memory safety},
    pdfnewwindow=true,      % links in new window
    colorlinks=true,        % false: boxed links; true: colored links
    linkcolor=black,        % color of internal links
    citecolor=black,        % color of links to bibliography
    filecolor=magenta,      % color of file links
    urlcolor=cyan           % color of external links
}

\begin{document}

% Generate title page.
\maketitle

% STEP 3:
% Write your abstract (of course you really do this last).
% You can make several abstract pages (if you want it in different languages),
% but you should also then redefine some of the above parameters in the proper
% language as well, in between the abstract definitions.

\begin{abstract}
\noindent Abstract to be written last.
\end{abstract}

% Place ToC
\mytableofcontents

\mynomenclature

% -----------------------------------------------------------------------------------
% STEP 4: Write the thesis.
% Your actual text starts here. You shouldn't mess with the code above the line except
% to change the parameters. Removing the abstract and ToC commands will mess up stuff.
\chapter{Introduction}

\begin{itemize}
  \item The 5G RAN control plane (the gNB-CU-CP) now runs as software on ordinary servers.
        It is large, it keeps state for every connected device, it serves many requests at
        once, and it has only a few milliseconds to complete each procedure~\cite{3gpp22261}.
        One defect, or one regression in the slowest requests, therefore degrades the whole
        site.
  \item This software is written mostly in C++. C++ gives the control over memory and timing
        that such a system needs, but it leaves memory-safety and concurrency mistakes to the
        programmer~\cite{szekeres2013sok}. Rust promises the same performance while rejecting
        those mistakes at compile time~\cite{klabnik2023rust}.
  \item The problem is a lack of evidence. Telecom vendors are under growing pressure to move
        to memory-safe languages~\cite{cisa2023memorysafe}, but no controlled study tells them
        what the change would mean for predictability, how many defects it would really
        prevent, or what the migration would cost for a representative control-plane
        subsystem~\cite{karlback2021ran}.
  \item Four research questions frame the work. RQ1 compares tail latency and jitter under
        overload~\cite{dean2013tail}. RQ2 asks which defect classes Rust removes at compile
        time. RQ3 asks what Rust's safety model costs in resources. RQ4 asks how expensive it
        is to adopt Rust step by step inside an existing C++ codebase.
  \item The work contributes four things: a benchmarking framework that others can rerun, one
        frozen specification implemented twice and tested for identical behaviour, a defect
        taxonomy that also marks where Rust's guarantee ends, and migration guidelines drawn
        from the measurements.
\end{itemize}

\chapter{Background: Systems Languages for Latency-Critical Software\label{chapter:background}}

\begin{itemize}
  \item In the NG-RAN architecture the base station (gNB) is split into a central and a
        distributed unit, and the central unit is split again into a control plane and a user
        plane. UE Context Management sits in the control plane, the gNB-CU-CP, and reaches its
        neighbours over the standardised F1, E1, and NG interfaces~\cite{3gpp38401}.
  \item Deadlines are met or missed because of predictability, not average speed. The metrics
        that matter are therefore the p99 and p99.9 latencies and the jitter around
        them~\cite{dean2013tail}. They have to be measured past the saturation point and
        recorded without coordinated omission, that is, without letting a load generator that
        has fallen behind hide its own slowest requests~\cite{tene_hdrhistogram}.
  \item C++ offers abstractions that cost nothing at run time, deterministic cleanup through
        RAII, and no garbage collector. Safety, however, remains the programmer's
        responsibility~\cite{szekeres2013sok}, and sanitizers report only the errors that a
        given test run happens to reach.
  \item Rust has the same run-time model and also no garbage collector, but it checks memory
        safety and freedom from data races at compile time, using ownership, borrowing, and the
        \texttt{Send}/\texttt{Sync} traits~\cite{jung2021safe}. Code that needs to step outside
        those rules must be marked \texttt{unsafe}, which keeps such code small and easy to
        find.
  \item The comparison is therefore C++ against Rust. Garbage-collected languages are left out
        because their pauses are incompatible with the latency targets. The concurrency model
        is held fixed in both implementations, and the C ABI is treated as the realistic
        adoption path. No earlier study isolates the language this way~\cite{karlback2021ran}.
\end{itemize}

\chapter{Problem Context: UE Context Management and Predictability\label{chapter:context}}

\begin{itemize}
  \item UE Context Management owns the lifecycle of one context per connected device. Each
        context moves through four states, Idle, Connected, Reconfiguring, and Released, driven
        by a small set of lifecycle operations. An operation that is illegal in the current
        state returns a typed error instead of being carried out.
  \item The subsystem was chosen because it is both representative and small enough to keep
        under control. It has mutable state per device, a strict lifecycle, heavy concurrency,
        and hard latency targets~\cite{3gpp22261}, and it is still simple enough to specify
        precisely and to test for identical behaviour in two languages.
  \item Five defect classes fall inside Rust's guarantee~\cite{jung2021safe}, and each has a
        natural form in this subsystem: use-after-free, double-free, iterator invalidation,
        data race, and unguarded shared mutation.
  \item Two further classes are included on purpose from outside that guarantee, a lock-order
        deadlock and a lost update. They keep the study from confirming a foregone conclusion
        and show instead where the guarantee stops.
  \item Scale is what makes the tail matter. One procedure triggers many operations and one
        site serves many devices, so a p99.9 measured on a single operation becomes a steady
        rate of degraded procedures in the field~\cite{dean2013tail}. This is why the load
        sweep is pushed past the point where the service starts to fall behind.
\end{itemize}

\chapter{Objectives\label{chapter:objectives}}

\begin{itemize}
  \item A documented literature search found no study that measures the tail latency of a Rust
        and a C++ implementation of the same stateful control-plane subsystem under
        overload~\cite{karlback2021ran}.
  \item Five objectives follow from the research questions: measure predictability (O1),
        establish which defect classes are prevented (O2), measure the resource overhead (O3),
        build and measure a migration prototype (O4), and publish reusable artifacts and
        guidance (O5).
  \item Three hypotheses are stated in advance, so that the measurements can refute them. H1:
        below saturation the two implementations perform the same~\cite{emmerich2019drivers}.
        H2: Rust prevents at compile time the defect classes for which C++ can only offer
        best-effort detection. H3: the cost of calling Rust from C++ is negligible.
  \item The work succeeds if the evidence is good enough to base a decision on: repeated runs
        reported with confidence intervals~\cite{georges2007rigorous}, defect claims shown on
        programs that really compile and run, migration cost read off a working prototype, and
        every number reproducible from the published artifacts.
  \item The scope is drawn narrowly to protect internal validity. The full gNB, the PHY and L2
        layers, the user plane, captured traffic traces, managed languages, and production
        operations are all left out.
\end{itemize}

\chapter{Research Methodology\label{chapter:methodology}}

\begin{itemize}
  \item The approach is empirical and constructive. Instead of comparing two systems that
        already exist and differ in many ways, the study builds a matched pair of
        implementations from one frozen specification, so that the language is the only
        systematic difference between them.
  \item RQ1 and RQ3 are answered with an open-loop load sweep that runs from light load into
        overload; open-loop means the generator keeps to its send schedule instead of waiting
        for replies. Latency is recorded in identically configured HDR
        histograms~\cite{tene_hdrhistogram}, and resource use is sampled through the operating
        system.
  \item RQ2 is answered by controlled fault injection~\cite{natella2016fault}. A fixed
        catalogue of defects is written in both languages, and the expected outcome of each one
        is recorded in advance, so that a surprise shows up as a discrepancy instead of being
        explained away afterwards.
  \item RQ4 is answered by a working prototype in which C++ drives the Rust core through a
        foreign-function interface. It yields three measurements: how large the seam is, how
        much \texttt{unsafe} code it needs, and what one call across it
        costs~\cite{li2022crosslang}.
  \item Fairness is the central experimental control: the specification is frozen, allocator
        and optimisation settings are matched, and every point where the two languages had to
        differ is written down in a forced-idiom log. Results are reported as medians with
        bootstrap confidence intervals, together with exact paired permutation tests corrected
        for multiple comparisons by the Holm procedure~\cite{georges2007rigorous}. All data is
        synthetic, so no data-protection concerns arise.
\end{itemize}

\chapter{Design and Implementation\label{chapter:design}}

\begin{itemize}
  \item A language-neutral specification is frozen before any code is written. It fixes the
        public API, the state machine, the data structures, the algorithmic strategy, and the
        sharding and concurrency model, and it serves as the fairness baseline for everything
        that follows.
  \item Both versions use the same shared-nothing, thread-per-shard architecture: eight shards,
        one bounded blocking MPSC queue per shard, single-use completion handles, strict FIFO
        order within a shard, and a single atomic counter as the only shared mutable state.
  \item The run-time substrate is equalised as well: the same jemalloc allocator, with the
        system allocator kept as a control, the same HDR histogram configuration, matching
        optimisation levels (\texttt{-O2} and \texttt{opt-level = 2}), and the same hash
        functions for choosing a shard and for lookups inside it.
  \item Where the languages must diverge, the forced-idiom log records why. C++ has to
        hand-write the bounded queue that Rust gets from \texttt{sync\_channel}, and the tagged
        request struct and response variant that C++ needs are larger and more branch-heavy than
        Rust's typed per-variant senders. The result is that the Rust version has more lines but
        lower total cyclomatic complexity~\cite{lizard}.
  \item The Rust core contains no \texttt{unsafe} at all. Equal behaviour is tested rather than
        assumed: a language-neutral conformance suite drives both implementations, and their
        canonical transcripts are compared line by line.
\end{itemize}

\chapter{Experimental Setup and Evaluation\label{chapter:eval}}

\begin{itemize}
  \item Two environments have to be kept apart. Everything reported here was measured on a WSL2
        development host, which is adequate for parity testing, fault injection, and the
        structure of the FFI seam, but not for claims about tail latency. The bare-metal
        protocol is specified and scripted, but it has not been run yet.
  \item The workload is a deterministic generator of RAN lifecycle traffic, justified by the
        specification, that follows a fixed open-loop send schedule and is swept from nominal
        load through saturation. The four variants, two languages by two allocators, are
        interleaved within every repeat.
  \item For RQ1 and RQ3, the four variants are indistinguishable at the median and equal in
        throughput and CPU use~\cite{emmerich2019drivers}. The tail cannot be resolved on this
        host, because the confidence interval on a single variant's p99 is wider than the spread
        between all four variants.
  \item For RQ3, memory is the one performance result that comes out cleanly. The Rust build
        carries a few MiB more peak resident set under either allocator, the difference has the
        same sign in every repeat, and it is larger than the difference the allocator makes.
  \item For RQ2, all five in-scope defect classes are compile errors in safe
        Rust~\cite{jung2021safe}, whereas every one of them compiles in C++ and is caught only
        if the right sanitiser happens to cover the path. The two boundary classes compile in
        both languages. For RQ4, a seam of twelve entry points lets C++ drive the Rust core,
        keeps that core free of \texttt{unsafe}, and costs less than a nanosecond per call.
\end{itemize}

\chapter{Discussion\label{chapter:discussion}}

\begin{itemize}
  \item Every research question is answered together with the status of its evidence. The safety
        and migration results are final, because they do not depend on the measurement
        environment. The performance and resource results are preliminary.
  \item On performance, the honest reading is that the question about the median is settled and
        the question about the tail is still open. That is the better position to be in: it
        leaves the eventual bare-metal measurement free to be a finding rather than a
        confirmation. The memory penalty is bounded, but explaining it would need hardware
        counters.
  \item On safety the difference is one of kind rather than degree, and it has a definite edge.
        Rust eliminates data races, not race conditions, and aliasing errors, not lock-ordering
        errors~\cite{jung2021safe}.
  \item On migration, an \texttt{unsafe}-free core says what the compiler checks; it says
        nothing about what a C++ host in the same address space can
        corrupt~\cite{li2022crosslang}. The gain is therefore that the component stops
        contributing new defects of its own.
  \item The hypotheses are then revisited: H1 is half answered, H2 holds in its precise form
        and fails in the loose reading that Rust prevents every concurrency bug, and H3 is
        supported. The chapter closes with engineering guidelines, related work, threats to
        validity, and the professional obligation not to over-claim.
\end{itemize}

\chapter{Conclusion and Future Work\label{chapter:conclusions}}

\begin{itemize}
  \item The thesis compared modern C++ and Rust along four dimensions, on one UE Context
        Management service built twice from a single frozen specification. Building both
        versions from one specification separates the effect of the language from differences
        in design or in the skill of the author.
  \item Safety: safe Rust rejects all five injected memory-safety and data-race classes at
        compile time, while the two boundary classes compile and misbehave exactly as they do
        in C++.
  \item Migration: an unmodified Rust core can be driven from C++ across a thin seam that
        confines \texttt{unsafe} to a few functions, at less than a nanosecond per call. The
        real cost of adoption is API design and the learning curve~\cite{fulton2021adopting}.
  \item Performance: the languages are equivalent at the median and in CPU use and throughput,
        the Rust build carries a few MiB more peak RSS, and the tail is not yet measurable on
        a shared host.
  \item Future work starts with the bare-metal campaign that is already specified, and then
        extends to production traces, further subsystems, hardware counters, async runtimes,
        in-process heap isolation, and a longitudinal migration case study.
\end{itemize}

% STEP 5:
% Uncomment the following lines and set your .bib file and desired bibliography style
% to make a bibliography with BibTeX.
% Alternatively you can use the thebibliography environment if you want to add all
% references by hand.

\cleardoublepage %fixes the position of bibliography in bookmarks
\phantomsection

\addcontentsline{toc}{chapter}{\bibname} % This lines adds the bibliography to the ToC
\bibliographystyle{IEEEtran} % numbered in citation order
\bibliography{bibliography}

\begin{appendices}
\myappendixtitle

\chapter{Fault-injection example\label{appendix:code}}

\begin{itemize}
  \item This appendix shows one case from the RQ2 catalogue, the use-after-free, in both
        languages side by side.
  \item In the C++ version a reference into the context map outlives the \texttt{erase} that
        frees the entry, and the program still compiles cleanly.
  \item The error appears only at run time, and only if AddressSanitizer is enabled on a path
        that a test actually exercises.
  \item The equivalent safe-Rust program does not compile at all: the borrow checker refuses to
        let the map be modified while a reference into it is still live
        (\texttt{E0499})~\cite{klabnik2023rust}.
  \item Side by side, the two programs make the difference in kind of guarantee concrete instead
        of arguing it from language semantics.
\end{itemize}

\end{appendices}

\end{document}
